\section{Conclusion}
\label{sec:conclusion}

Evaluated under realistic conditions---chronological replay, attack classes that appear after training, and a
small, delayed labelling budget---intrusion detectors that are trained once collapse, and simple retraining
schemes forget or flood analysts with reviews. \nova{} combines class-conditional open-world rejection, a
novelty-weighted and diverse query, conservative label propagation, a class-balanced memory and a guarded
deployment loop. At a 10\% budget it is the best or tied-best method on open-world macro-F1 on all three
datasets and has the lowest false-positive rate on all three. At a 1\% budget it leads on both CIC datasets
(by 9.0 and 3.4 points). On UNSW-NB15 at 1\% it trails methods that never learn the new classes and keep
rejecting them, while still halving the false-positive rate. Its
accuracy lead did not carry over to an MLP classifier, and short-lived attacks that look like benign traffic
remain undetected by every method we tried, including a meta-learned embedding. Both are open problems for
future work. All runs,
configurations and per-flow predictions are released with the SentinelAI console.
